% D.81 candidate. UNREGISTERED, NOTES ONLY, pending full independent audits.
\subsection{Every independent completed unit pair, including nonprimitive local lines}

\begin{statement}{D.81}{Open}{Candidate: the full compact kernel for every independent unit pair}
Let \(a,b\in H^1(G_{\mathbb Q},\mathbb Z_5(1))\) have zero
valuations at five and linearly independent reductions modulo five.
Their ramification may be infinite. Write their local unit classes as
\[
 a|_{G_{\mathbb Q_5}}=\alpha\kappa(6),\qquad
 b|_{G_{\mathbb Q_5}}=\beta\kappa(6),
 \qquad \alpha,\beta\in\mathbb Z_5.
                                                               \tag{D.81.1}
\]
There is no primitivity hypothesis on this local vector. Define
\[
 \begin{split}
 \epsilon(\alpha,\beta)&=
 \begin{cases}
  1,&(\bar\alpha,\bar\beta)=(0,0)
       \text{ or }\bar\alpha\bar\beta(\bar\alpha-\bar\beta)\ne0,\\
  0,&\text{otherwise},
 \end{cases}\\
 n(\alpha,\beta)&=
      v_5(\alpha)+v_5(\beta)+v_5(\alpha-\beta),
       \qquad v_5(0)=\infty.
 \end{split}                                                   \tag{D.81.2}
\]
Use \(Z(n)=\mathbb Z_5^n\) at finite \(n\), and
\(Z(\infty)=\prod_{\mathbb N}\mathbb Z_5\), as in D.77.
For the actual Alexander module \(A_s\), \(s=(a,b)\), the candidate gives
\[
 \mathcal S_s\simeq(\mathbb Z/5)^{\epsilon(\alpha,\beta)},
 \qquad \mathcal T_s=0,\qquad
 \mathcal K_s\simeq
       (\mathbb Z/5)^{\epsilon(\alpha,\beta)}
                         \times Z(n(\alpha,\beta)).
                                                               \tag{D.81.3}
\]
The last isomorphism is continuous and noncanonical. The finite term
is the inverse finite-global-localization kernel of D.63, not merely
a restriction kernel with divisible local coefficients.
There are thirteen residual vectors giving the order-five term:
the zero vector and the twelve vectors off the three boundary lines.
The twelve nonzero boundary vectors give no such term.

In particular, explicit finitely ramified actual Alexander modules
have compact kernels
\[
  \mathbb Z/5\times\mathbb Z_5^n\quad(n\ge3),
  \qquad
  \mathbb Z/5\times\prod_{\mathbb N}\mathbb Z_5.
                                                               \tag{D.81.4}
\]
Their completed Kummer parameters are given in (D.81.20) below.
These examples are assertions about actual Galois modules and their
cohomology. They do not assert that the corresponding abelian pairs
have metabelian sections. General first-cohomology ranks with
infinite ramification, the general nonseparated kernel, and the
first higher-derived lifting obstructions remain outside this
classification.
\end{statement}

\begin{proof}[Candidate proof, pending independent audits]
\paragraph{The actual images do not require a primitive local line.}
Let \(\Lambda=\mathbb Z_5[[u,v]]\), \(t_1=1+u,t_2=1+v\), and let
\(I_0\) be the actual untwisted operator image in D.78.
The two global hypotheses give two residual inertia vectors away
from five with unit determinant, exactly as in D.80. Their pure
translation images generate \(H=\mathbb Z_5^2\). Hence
\[
                         I_s=H\rtimes I_0.             \tag{D.81.5}
\]
The convergence and exact Kummer interpretation of (D.81.1) are
also those proved in D.80. Apply its integral local monomial gauge
so that (D.81.1) holds for actual cocycles. D.78 proves both
surjectivity of \(G_{\mathbb Q_5}\to I_0\) and
\(\kappa(6)(\ker(G_{\mathbb Q_5}\to I_0))=\mathbb Z_5\).
Consequently, with its actual zero-translation complement,
\[
 I_5=H_5\rtimes I_0,\qquad
 H_5=\mathbb Z_5(\alpha,\beta).                         \tag{D.81.6}
\]
This holds when the line is nonprimitive and when \(H_5=0\).
Indeed every translation coordinate lies in \(H_5\), and the
kernel just described realizes every element of that subgroup,
so a lift of any \(I_0\)-element can be corrected to translation zero.
The gauge preserves (D.81.5), the constant dual submodule, and
all finite augmentation quotients.

Put \(D=A_s^D\), with the Cartier-dual convention of D.78.
Translation coinduction gives
\[
 H^j(H,D)=H^j(H_5,D)=0\quad(j>0).                       \tag{D.81.7}
\]
For the second assertion no saturated complement is presumed.
If \(H_5\ne0\), an integral linear change of coordinates identifies
it with \(5^e\mathbb Z_5\times0\). A continuous transversal is
\(\{0,\ldots,5^e-1\}\times\mathbb Z_5\); thus as a compact
translation space \(H\) is \(H_5\) times that transversal.
The contraction of continuous function-valued cochains in D.78
works with the transversal as an additional coefficient variable.
When \(H_5=0\), the positive-degree assertion is immediate.

Set
\[
 \begin{gathered}
 f=t_1^\alpha t_2^\beta-1,\qquad R=\Lambda/(f),\qquad
 J=(u,v)R,\\
 D_{\rm lin}=R^D=D^{H_5},\qquad
 \lambda=(R/J)^D=\mathbb Q_5/\mathbb Z_5(-1).
 \end{gathered}                                               \tag{D.81.8}
\]
If \(\alpha=\beta=0\), this means \(R=\Lambda\).
The augmentation character on \(R/J\) is \(\chi^2\), which
explains the single inverse cyclotomic character after Cartier duality.
Continuous Hochschild--Serre now identifies restriction with
\[
 H^1(I_0,\lambda)\longrightarrow H^1(I_0,D_{\rm lin}).
                                                               \tag{D.81.9}
\]
The finite affine-image Chebotarev argument of D.78/D.80 applies
without bounded ramification and gives the injection
\[
 \mathcal S_s^*\lhook\joinrel\longrightarrow
 \ker\{H^1(I_0,\lambda)\to H^1(I_0,D_{\rm lin})\}.
                                                               \tag{D.81.10}
\]
Here the star is ordinary Pontryagin duality. That argument is
at a finite coefficient level before taking its direct limit:
a locally trivial finite cocycle is zero on the kernel of its
actual finite operator image, because any nonzero affine value
there would occur at an unramified Frobenius. Thus no unrestricted
coefficient action or formal coefficient assignment is substituted
for the actual image in (D.81.10).

\paragraph{A connecting module valid even at zero or nonprimitive slope.}
Write \(w=t_1t_2-1\), \(P=uvw\), and
\(N_0=\ker(I_0\to\Gamma)\), where \(\Gamma=\mathbb Z_5^\times\).
For every \(\tau\in N_0\), the exact three boundary formulas in
D.39 give
\(\mathcal B_\tau-1\in P\Lambda\).
D.77 supplies an actual local element \(\delta\) for which
\[
              \mathcal B_\delta-1=P E,\qquad
                                  E\in\Lambda^\times. \tag{D.81.11}
\]
Its image in every \(R\) remains a unit. Since the \(N_0\)-action
is multiplication by the corresponding beta series, it follows
on the two actual dual modules that
\[
 \begin{split}
 D_{\rm lin}^{I_0}&=((R/PR)^D)^\Gamma,\\
 (D_{\rm lin}/\lambda)^{I_0}&=((J/PJ)^D)^\Gamma.
 \end{split}                                                 \tag{D.81.12}
\]
The second identity uses \(D_{\rm lin}/\lambda=J^D\).
On these quotient modules the action factors through \(\Gamma\);
no finite-rank hypothesis on either quotient is needed.

The actual ideal inclusion gives the compact exact sequence
\[
 0\longrightarrow C\longrightarrow J/PJ
       \longrightarrow R/PR\longrightarrow\mathbb Z_5(2)
       \longrightarrow0,\qquad
 C=PR/PJ=\frac{R}{J+\operatorname{Ann}_R(P)}.
                                                               \tag{D.81.13}
\]
We compute its first term explicitly, including every zero-divisor case.
The quotients \(\Lambda/(u),\Lambda/(v),\Lambda/(w)\) are domains.
For one of these boundary elements \(L\), if \(f\bmod L\ne0\),
then \(L\) is a nonzero divisor in \(R\): from \(Lr=fh\) in
\(\Lambda\), reduction modulo \(L\) forces \(h=Lh_1\);
cancellation gives \(r=fh_1\).

For a nonzero vector \((\alpha,\beta)\), the three reductions
of \(f\) vanish respectively exactly when
\(\beta=0\), \(\alpha=0\), and \(\alpha=\beta\).
At most one of these equalities occurs. If none occurs, all
three factors of \(P\) are nonzero divisors, so
\(\operatorname{Ann}_R(P)=0\). If one occurs, write
\(f=L S_\eta(L)\), where
\[
 S_\eta(L)=\frac{(1+L)^\eta-1}{L},\qquad S_\eta(0)=\eta\ne0.
 \]
Here \(\eta=\alpha\), \(\beta\), or their common value,
according to the boundary. The other two factors of \(P\)
are nonzero divisors in \(R\); cancellation in \(\Lambda\)
then gives
\(\operatorname{Ann}_R(P)=\operatorname{Ann}_R(L)=(S_\eta(L))\).
This argument also applies when \(\eta\) is a unit.
If \(\alpha=\beta=0\), the original ring \(\Lambda\) is a
domain and \(P\) is again a nonzero divisor. Thus
\[
 C\simeq
 \begin{cases}
  \mathbb Z_5(5),&
       \alpha\beta(\alpha-\beta)\ne0
                     \text{ or }(\alpha,\beta)=(0,0),\\
  (\mathbb Z_5/\eta\mathbb Z_5)(5),&
       \text{a nonzero exact boundary with exponent }\eta .
 \end{cases}                                                 \tag{D.81.14}
\]
To verify the indicated arithmetic character, the multiplier
of \(\rho_0(g)\) has constant term \(\chi(g)^2\), and
\(\varphi_\chi(P)/P\) has constant term \(\chi(g)^3\).
The action on the generator \(P\bmod PJ\) is therefore
\(\chi(g)^5\), including in the finite quotient case.
No division by \(\eta\), or by its five-adic valuation, occurs.

Cartier duality of (D.81.13) is exact and gives
\[
 0\longrightarrow\lambda\longrightarrow (R/PR)^D
       \longrightarrow (J/PJ)^D\longrightarrow C^D
       \longrightarrow0.                                \tag{D.81.15}
\]
The order-four subgroup of \(\Gamma\) acts nontrivially
on \(\lambda\), with unit difference from one. Averaging
and Hochschild--Serre give
\(H^0(\Gamma,\lambda)=H^1(\Gamma,\lambda)=0\).
Taking the two short exact sequences in (D.81.15), and
using (D.81.12), therefore injects the entire restriction
kernel in (D.81.10) into \((C^D)^\Gamma\).
In the infinite case this last group is
\((\mathbb Q_5/\mathbb Z_5(-4))^\Gamma=\mathbb F_5\).
In the finite case it is zero for \(v_5(\eta)=0\), and
has order five for \(v_5(\eta)\ge1\).
Indeed \(2\) topologically generates \(\mathbb Z_5^\times\),
and \(v_5(2^4-1)=1\). In particular
\[
                         |\mathcal S_s|\le5            \tag{D.81.16}
\]
for every pair in the statement. This bounds the full kernel,
not only its elements of order five.

\paragraph{One fixed primitive for every locally five-divisible pair.}
When \((\alpha,\beta)\notin5\mathbb Z_5^2\), D.80 gives
precisely (D.81.3). It remains to prove the nonzero lower bound
when \(\alpha,\beta\in5\mathbb Z_5\).

Let \(c=48h\bmod5\) be the actual integral weight-three cocycle
of D.31, as in D.78. Its coefficients may be viewed in
\(\lambda[5]\), because \(\chi^3=\chi^{-1}\) modulo five.
Its class in \(H^1(I_0,\lambda)\) is nonzero.
Indeed its restriction to \(N_0\) is primitive raw
\(\kappa_3\bmod5\), so its class with \(\lambda[5]\)-coefficients
is nonzero. The map from that cohomology group to
\(H^1(I_0,\lambda)\) is injective since \(\lambda^{I_0}=0\).
Actual coinduction and inflation therefore keep this class
nonzero in \(H^1(G_{\mathbb Q},D)\).

Use the one finite module
\[
 M=(A_s/(5,(u,v)^4)A_s)^D.
                                                               \tag{D.81.17}
\]
The exact local gauge in (D.81.1) makes every local translation
exponent divisible by five. On \(M\) these translations are
the identity, since \(t_i^{5z}=1\) modulo \((5,(u,v)^4)\).
The whole local action on \(M\) is therefore the untwisted one.

Let \(\varepsilon_{ij}\) extract the coefficient of \(u^iv^j\)
in the compact ten-dimensional quotient, with \(i+j\le3\).
The single functional
\[
 m=2\varepsilon_{01}+4\varepsilon_{10}
         +2\varepsilon_{11}+\varepsilon_{20}
         +2\varepsilon_{21}\quad\text{over }\mathbb F_5
                                                               \tag{D.81.18}
\]
satisfies \((g-1)m=c(g)\varepsilon_{00}\) for every actual
untwisted action through this quotient. Here is a finite
coefficient verification, retaining the actual arithmetic
normalization. Its compact operator is
\[
 \begin{split}
 \rho_0(g)F&=S_\chi(t_1)S_\chi(t_2)\mathcal B_g
                         F(t_1^\chi-1,t_2^\chi-1),\\
 \mathcal B_g&\equiv
  1+\frac{\chi^2-1}{24}uv
   +\left(\frac{c(g)}{48}-\frac{\chi^2-1}{48}\right)
                                  (u^2v+uv^2)
                 \pmod{5,(u,v)^4}.
 \end{split}                                                 \tag{D.81.19}
\]
This is exactly the D.31/D.39 beta series; all denominators
are units at five. The coefficients of \(S_\chi\) and of
the substitutions through degree three are binomial
polynomials with unit denominators. On the ten monomials,
expansion gives
\(\chi m-\rho_0(g)^t m=c(g)\chi^2\varepsilon_{00}\)
over \(\mathbb F_5[\chi,c]/(\chi^4-1)\).
More explicitly, before imposing \(\chi^4=1\), the ten
residuals over \(\mathbb F_5[\chi,c]\), in the monomial
order \(1,v,u,v^2,uv,u^2,v^3,uv^2,u^2v,u^3\), are
\[
 \chi(\chi^4-1)
             (0,-2,1,0,-2,-1,0,0,-2,0).
\]
Cartier action is \(\chi\rho_0(g)^{-t}\), so this identity
is precisely the asserted coboundary equation.
The accompanying certificate verifies the ten symbolic
remainders and all twenty residual pairs \((\chi,c)\);
it does not assign any unknown higher beta coefficient.
Thus a single primitive in \(M\) works for the entire
local decomposition group at five.

The finite-level uniformity at all other primes is now
exactly the argument proved in D.80. The set
\(E=\{q\ne5:(a_q,b_q)\notin5\mathbb Z_5^2\}\) is finite.
At every prime outside \(E\cup\{5\}\), \(M\) is unramified
even if the full module is ramified, and D.78's universal
Frobenius primitives apply. At each prime in \(E\),
D.43 and finite local Tate duality, with the finite
Mittag--Leffler comparison in D.80, give
\(H^1(G_{\mathbb Q_q},D)=0\). Its one local class therefore
dies at some larger finite coefficient level.
Taking the maximum of these finitely many levels and
the level containing \(M\) gives one common finite
global coefficient module. The primitive (D.81.18)
handles five in that same module; real cohomology is zero.
The nonzero global class \(c\) consequently belongs to
the direct limit of finite everywhere locally trivial
classes. Equation (D.81.16) gives
\(\mathcal S_s=\mathbb Z/5\).
No direct limit has been interchanged with an infinite
product of local cohomology groups.

Global independence modulo five makes \(a,b,a-b\)
all primitive. D.77 therefore gives
\(\mathcal T_s=0\) and the continuous decomposition in
(D.81.3), with exactly the three local valuations in
(D.81.2). Counting the residual vectors gives twelve
off-boundary vectors, one zero vector, and twelve
nonzero boundary vectors, as stated.

\paragraph{Explicit mixed finite and free compact kernels.}
For \(q\ne5\), put
\(\lambda_q=\log(q^4)/(4\log6)\in\mathbb Z_5\).
We have \(\lambda_{31}\equiv1\pmod5\):
\(31^4\equiv21\pmod{25}\), and \(4\log6\equiv20\pmod{25}\).
For any prescribed \(A,B\in\mathbb Z_5\), define actual
completed Kummer classes
\[
 \begin{split}
 a&=\kappa(7)+
       \frac{A-\lambda_7}{\lambda_{31}}\kappa(31),\\
 b&=\kappa(43)+
       \frac{B-\lambda_{43}}{\lambda_{31}}\kappa(31).
 \end{split}                                                 \tag{D.81.20}
\]
All coefficients are integral. These classes are supported
on \(\{7,31,43\}\), have zero five-valuations, and their
inertia vectors at seven and forty-three are the standard
basis. Their local logarithms are exactly \(A,B\).
For each integer \(n\ge3\), take
\[
                 A=5,\qquad B=5(1+5^{n-3}).
                                                               \tag{D.81.21}
\]
Their two valuations are one, and
\(v_5(A-B)=n-2\). Their residual vector is zero.
Equation (D.81.3) gives the first family in (D.81.4).
Taking \(A=B=0\) gives its countable-product example.
These are actual finitely ramified Galois modules, so
D.74 also computes their unrestricted first cohomology
as \(\prod_{\mathbb N}\mathbb Z_5\).
The existence of a section above an arbitrary one of
these completed abelian pairs is not inferred from
either cohomology computation.
\end{proof}
